\documentclass[11pt,a4paper]{article}
\usepackage{gtu-el-yazisi}
\begin{document}
\notbilgi{Alan Teorisi}{Uygulama 8 -- İletkenler: potansiyel, yüzey yükü, kapasite}{Uygulama-8.pdf}
% yerel yardımcı: üstü çizili (sadeleştirilen) terimler
\newcommand\iptalX[3]{\tikz[baseline=(c.base)]{\node[inner sep=0pt](c){$#2#3$};\draw[#1,line width=0.5pt](c.south west)--(c.north east);}}
\newcommand\iptalK[2]{\iptalX{murekkep}{#1}{#2}}
\newcommand\iptalO[2]{\iptalX{turuncukalem}{#1}{#2}}
\newcommand\iptalR[2]{\iptalX{kirmizikalem}{#1}{#2}}
\newcommand\iptal[1]{\mathpalette\iptalK{#1}}
\newcommand\iptalt[1]{\mathpalette\iptalO{#1}}
\newcommand\iptalk[1]{\mathpalette\iptalR{#1}}

\sayfa{1}
\noindent\textbf{a)}
\begin{center}
\begin{tikzpicture}[line width=0.8pt]
  \draw[-{Stealth}] (0,-1.1) -- (0,3.3);
  \draw[-{Stealth}] (-2.0,2) -- (6.6,2) node[right] {$V_d$};
  \draw[-{Stealth}] (-2.0,0) -- (6.6,0) node[right] {$V_0$};
  \foreach \hx in {-1.8,-1.3,...,6.2} {\draw[line width=0.6pt] (\hx,2) -- ++(-0.2,0.22); \draw[line width=0.6pt] (\hx,0) -- ++(-0.2,-0.22);}
  \node at (0.3,2.4) {$d$};
  \node[left] at (-0.05,0.3) {$O$};
  \node at (2.4,2.7) {$\varepsilon_3$\ (3)};
  \node at (1.6,1.0) {$\varepsilon$};
  \node at (2.6,1.0) {(2)};
  \node at (1.6,-0.75) {$\varepsilon_1$};
  \node at (2.6,-0.75) {(1)};
  \draw[kirmizikalem,-{Stealth}] (6.0,2.2) -- (6.0,2.9);
  \node[kirmizikalem,right] at (6.0,2.7) {$\vec n=\vec e_x$};
\end{tikzpicture}
\end{center}

Ara bölgede potansiyel fonksiyonunu, iletkenler üzerinde birim alanda toplanan yükleri bulunuz. $x=d$'de birim alanda toplanan yük $Q_d$ ise $C=Q_d/(V_d-V_0)$'ı bulunuz.

\bigskip
\noindent Ara bölgede \quad $\Delta V=0$
\[ \frac{\partial^2 V}{\partial x^2}=0 \qquad\qquad V(x)=Ax+B \]
\begin{align*}
x=0\ &\Rightarrow\ V_0=A\cdot 0+B \\
 &\qquad B=V_0 \\
x=d\ &\Rightarrow\ V_d=Ad+V_0 \\
 &\qquad A=\frac{V_d-V_0}{d}
\end{align*}
\[ \text{O halde}\qquad \kutula{V(x)=\frac{V_d-V_0}{d}\,x+V_0} \]
\[ \vec E=-\mathrm{grad}\,V=-\frac{V_d-V_0}{d}\,\vec e_x \]
\noindent $x=d$ düzleminde
\begin{align*}
\varepsilon_3\,\overset{\textcolor{kirmizikalem}{0}}{\iptalk{E_n^{(3)}}}-\varepsilon E_n &= \rho_{s,d} \qquad\qquad \textcolor{kirmizikalem}{\shortstack[l]{iletken ortamda\\ Elektrik alan 0}} \\
\rho_{s,d} &= \frac{\varepsilon}{d}\,(V_d-V_0) \\
Q_d &= \int_S \rho_{s,d}\,ds=\frac{\varepsilon}{d}\,(V_d-V_0) \qquad\qquad \textcolor{kirmizikalem}{\shortstack[l]{iletkenler kapattığı bölgenin\\ içinde olan elektrik alanı\\ dışarıya geçirmez.}} \\
C &= \frac{Q_d}{V_d-V_0}=\frac{\varepsilon}{d}
\end{align*}
\noindent $x=0$ düzleminde
\[ \varepsilon E_n^{(2)}-\varepsilon_1\,\overset{\textcolor{kirmizikalem}{0}}{\iptalk{E_n^{(1)}}}=\rho_{s,0} \]
\[ \rho_{s,0}=-\frac{\varepsilon}{d}\,(V_d-V_0)=Q_0 \qquad \text{(Birim alanda)} \]

\sayfa{2}
\noindent\textbf{b)}
\begin{center}
\begin{tikzpicture}[line width=0.8pt]
  \draw[-{Stealth}] (0,-1.4) -- (0,3.9) node[right] {$x$};
  \draw[-{Stealth}] (-2.3,2.7) -- (6.0,2.7) node[right] {$V_d$};
  \draw[-{Stealth}] (-2.3,0) -- (6.0,0) node[right] {$V_0$};
  \draw[dashed] (-2.3,1.35) -- (5.5,1.35);
  \foreach \hx in {-2.0,-1.6,...,5.4} \draw[line width=0.6pt] (\hx,0) -- ++(-0.2,-0.25);
  \node at (-0.4,2.45) {$d$};
  \node at (-0.4,1.4) {$c$};
  \node at (-0.4,-0.6) {$O$};
  \node at (2.9,3.3) {(4)};
  \node at (4.3,3.3) {$\varepsilon$};
  \node at (2.0,1.95) {$\varepsilon_2$};
  \node at (3.2,1.95) {(3)};
  \node at (2.0,0.6) {$\varepsilon_1$};
  \node at (3.2,0.6) {(2)};
  \node at (2.9,-0.8) {(1)};
  \node at (4.5,-0.8) {$\varepsilon$};
\end{tikzpicture}
\end{center}

\[ \text{Ara bölgede}\quad \Delta V=0 \qquad\qquad V(x)=\begin{cases} Ax+B\,, & 0<x<c \\ Dx+F\,, & c<x<d \end{cases} \]
\noindent
\begin{minipage}[c]{0.28\linewidth}
\centering
$\boxed{\begin{array}{l} x=0\\ V(0)=V_0\\ B=V_0 \end{array}}$
\end{minipage}\hfill
\begin{minipage}[c]{0.68\linewidth}
\begin{align*}
V(d) &= V_d \\
Dd+F &= V_d\ \Rightarrow\ F=V_d-Dd
\end{align*}
\end{minipage}
\begin{align*}
x=c\qquad V^{(2)}(c) &= V^{(3)}(c) \\
Ac+V_0 &= Dc+F \\
Ac+V_0 &= Dc+V_d-Dd
\end{align*}
\[ \kutula{Ac+V_0=D(c-d)+V_d} \]
\noindent $x=c$
\begin{align*}
\varepsilon_2 E_n^{(3)}-\varepsilon_1E_n^{(2)} &= 0 \qquad \leftarrow\ \text{yük yok} \\
\varepsilon_1\frac{\partial V^{(2)}}{\partial x}-\varepsilon_2\frac{\partial V^{(3)}}{\partial x} &= 0
\end{align*}
\[ \kutula{\varepsilon_1A=\varepsilon_2D\ \Rightarrow\ D=A\,\frac{\varepsilon_1}{\varepsilon_2}} \]
\[ Ac+V_0=A\,\frac{\varepsilon_1}{\varepsilon_2}\,(c-d)+V_d \]
\[ A=\frac{\varepsilon_2(V_d-V_0)}{\varepsilon_1 d+c(\varepsilon_2-\varepsilon_1)} \]
\[ D=\frac{\varepsilon_1(V_d-V_0)}{\varepsilon_1 d+c(\varepsilon_2-\varepsilon_1)} \]

\sayfa{3}
\[ \kutula{F=V_d-Dd=\frac{c\,V_d(\varepsilon_2-\varepsilon_1)+\varepsilon_1 d\,V_0}{\varepsilon_1 d+c(\varepsilon_2-\varepsilon_1)}} \]
\medskip\hrule\medskip
\[ V(x)=\begin{cases}
\dfrac{\varepsilon_2(V_d-V_0)}{\varepsilon_1 d+c(\varepsilon_2-\varepsilon_1)}\,x+V_0\,, & 0<x<c \\[4mm]
\dfrac{\varepsilon_1(V_d-V_0)}{\varepsilon_1 d+c(\varepsilon_2-\varepsilon_1)}\,x+\dfrac{c\,V_d(\varepsilon_2-\varepsilon_1)+\varepsilon_1 d\,V_0}{\varepsilon_1 d+c(\varepsilon_2-\varepsilon_1)}\,, & c<x<d
\end{cases} \]
\[ \vec E(x)=\begin{cases}
-\dfrac{\varepsilon_2(V_d-V_0)}{\varepsilon_1 d+c(\varepsilon_2-\varepsilon_1)}\,\vec e_x\,, & 0<x<c \\[4mm]
-\dfrac{\varepsilon_1(V_d-V_0)}{\varepsilon_1 d+c(\varepsilon_2-\varepsilon_1)}\,\vec e_x\,, & c<x<d
\end{cases} \]
\noindent\textcolor{kirmizikalem}{$x=0$ ve $x=d$ deki yük yoğunlukları}
\begin{align*}
x=0:\quad \rho_{s,0} &= \varepsilon_1 E_n^{(2)}-\varepsilon\,\overset{\textcolor{kirmizikalem}{0}}{\iptalk{E_n^{(1)}}} \\
\rho_{s,0} &= -\frac{\varepsilon_1\varepsilon_2(V_d-V_0)}{\varepsilon_1 d+c(\varepsilon_2-\varepsilon_1)}
\end{align*}
\begin{align*}
x=d:\quad \rho_{s,d} &= \varepsilon E_n^{(4)}-\varepsilon_2 E_n^{(3)} \\
\rho_{s,d} &= \frac{\varepsilon_1\varepsilon_2(V_d-V_0)}{\varepsilon_1 d+c(\varepsilon_2-\varepsilon_1)}
\end{align*}

\bigskip
\noindent\textcolor{kirmizikalem}{\underline{S50 Pr.5}}
\begin{center}
\begin{tikzpicture}[line width=0.8pt,scale=0.8]
  \draw[smooth cycle] plot coordinates {(-1.9,1.5) (-0.4,1.9) (1.4,1.5) (2.5,0.1) (2.0,-1.4) (0.3,-2.1) (-1.4,-1.8) (-2.7,-0.9) (-2.5,0.5)};
  \draw[-{Stealth}] (0.25,1.85) -- (0.25,2.6);
  \node[right] at (0.25,2.45) {$\hat n$};
  \node at (-1.5,2.1) {(2)};
  \node at (-1.15,0.9) {$S$};
  \node at (0.5,0.9) {(1)};
  \node at (-1.3,-0.1) {iletken};
\end{tikzpicture}
\end{center}

Cismin dışında $V$ potansiyeli biliniyor olsun. Yüzeyinde toplanan yük:
\[ Q=-\int_S \varepsilon\,\frac{\partial V}{\partial n}\,ds \]
\noindent $S$ yüzeyinde
\begin{align*}
\varepsilon E_n^{(2)}-\varepsilon_1\,\overset{\textcolor{kirmizikalem}{0}}{\iptalk{E_n^{(1)}}} &= \rho_s \qquad\qquad \textcolor{kirmizikalem}{\shortstack[l]{iletkenlerin sadece\\ yüzeyinde yük birikir}} \\
\rho_s &= \varepsilon E_n^{(2)} \\
\rho_s &= -\varepsilon\,\frac{\partial V}{\partial n}
\end{align*}

\sayfa{4}
\[ Q=\int_S\rho_s\,ds=-\int_S\varepsilon\,\frac{\partial V}{\partial n}\,ds \]

\noindent\textcolor{kirmizikalem}{\underline{S.50 Pr.6}}
\begin{center}
\begin{tikzpicture}[line width=0.8pt,scale=0.9]
  \draw[-{Stealth}] (0,0) -- (0,3.4) node[right] {$z$};
  \draw[-{Stealth}] (0,0) -- (3.9,0) node[right] {$y$};
  \draw[-{Stealth}] (0,0) -- (-3.05,-1.85) node[below] {$x$};
  \draw[pattern=north east lines] (0,0) circle (1.1);
  \draw (0,0) -- (-15:1.1);
  \node[right] at (1.1,-0.35) {$R$};
  \node[right] at (0.6,-1.2) {iletken};
  \draw[-{Stealth}] (0.6,-1.15) -- (0.35,-0.8);
\end{tikzpicture}
\end{center}

\noindent küre üzerinde potansiyel $V_0$'dır\\
küre üzerindeki $Q$ yükünü\\
ve $C=\dfrac{Q}{V_0}$'ı bulunuz
\[ \Delta V=0 \]
\[ \frac{1}{r^2}\,\frac{\partial}{\partial r}\Bigl(r^2\frac{\partial V}{\partial r}\Bigr)=0\ \Rightarrow\ V(r)=\frac{A}{r}+B \]
\[ r\to\infty,\ V(r)\to 0\ \Rightarrow\ B=0 \]
\[ r=R,\ V(R)=V_0\ \Rightarrow\ \frac{A}{R}=V_0\ \Rightarrow\ A=V_0R \]
\[ V(r)=\frac{V_0R}{r} \]
\[ Q=-\int_S \varepsilon\,\frac{\partial V}{\partial n}\,ds=+\int_{\theta=0}^{\pi}\int_{\theta=0}^{2\pi}\varepsilon\,\frac{V_0R}{\iptal{r^2}}\,\underbrace{\iptal{r^2}\sin\theta\,d\theta\,d\phi}_{ds}\ \footnotemark \]
\footnotetext{Claude notu: muhtemelen ikinci integralin alt sınırı $\phi=0$ olmalı.}
\[ Q=4\pi\varepsilon V_0R \]
\[ C=\frac{Q}{V_0}=4\pi\varepsilon R \]

\bigskip
\begin{minipage}{\linewidth}
\noindent\textcolor{kirmizikalem}{\underline{S.50 Pr.7}}
\begin{center}
\begin{tikzpicture}[line width=0.8pt]
  \draw (-4.0,0.9) -- (0,0.9);
  \draw (-4.0,-0.9) -- (0,-0.9);
  \draw (0,0) ellipse (0.5 and 0.9);
  \draw[-{Stealth}] (-4.5,0) -- (2.6,0) node[below right] {$z$};
  \draw (0,0) -- (0,0.9);
  \node[right] at (0.05,0.45) {$R$};
  \draw[dashed] (0,0) ellipse (1.1 and 1.6);
  \node at (0.6,1.85) {$R=1$};
  \node at (1.5,1.2) {$V_0$};
  \node at (0.45,-1.0) {$V_R$};
\end{tikzpicture}
\end{center}
\end{minipage}

\sayfa{5}
\noindent Silindirin birim uzunluğundaki $Q$ yükünü\\
ve $C=\dfrac{Q}{V_0}$'ı bulunuz

\medskip
\noindent $\rho>R$
\begin{align*}
\Delta V=0\qquad & V(\rho)=A\ln\rho+B \\
\rho=1,\qquad & V_0=A\ln 1+B\ \Rightarrow\ B=V_0 \\
\rho=R,\qquad & V_R=A\ln R+V_0\ \Rightarrow\ A=\frac{V_R-V_0}{\ln R}
\end{align*}
\[ V(\rho)=\frac{\ln(\rho)}{\ln R}\,(V_R-V_0)+V_0 \]
\[ Q=-\int_S\varepsilon\,\frac{\partial V}{\partial n}\,ds=-\varepsilon\int_{z=0}^{1}\int_{\phi=0}^{2\pi}\frac{V_R-V_0}{\rho\ln R}\,\iptal{\rho}\,d\phi\,dz \]
\[ Q=-2\pi\varepsilon\,\frac{V_R-V_0}{\ln R} \qquad\Rightarrow\qquad C=\frac{Q}{V_0}=-\frac{2\pi\varepsilon\,(V_R-V_0)}{V_0\ln R} \]

\bigskip
\noindent\textcolor{kirmizikalem}{\underline{S.50 Pr.8}}

\elcizimi[0.6]{sekiller/uygulama-notlari-Uygulama-8-s5-1.png}[5]

\noindent $a$ yarıçaplı kürecik ve çeper iletkendir

\medskip
\noindent a)\ Oyuk içindeki Elektrik alanı bulunuz
\[ \textcolor{kirmizikalem}{\int_S \vec D\cdot d\vec s=Q} \]
\[ \textcolor{kirmizikalem}{\int_{\theta=0}^{\pi}\int_{\phi=0}^{2\pi}\varepsilon E_r\,r^2\sin\theta\,d\phi\,d\theta=Q} \]
\[ \vec E_r=\frac{Q}{4\pi\varepsilon r^2}\ \Rightarrow\ \vec E=\frac{Q}{4\pi\varepsilon r^2}\,\vec e_r \]

\noindent b)\ kabuğun iç ve dış yüzeylerinde biriken yükler nedir?
\begin{align*}
r=b\qquad \rho_{s,b} &= \overset{0}{\iptal{\varepsilon_1E_n^{(2)}}}-\varepsilon E_n^{(1)} \\
\rho_{s,b} &= -\frac{\varepsilon Q}{4\pi\varepsilon b^2}=-\frac{Q}{4\pi b^2}
\end{align*}

\sayfa{6}
\[ \rho_{s,c}=-\frac{\varepsilon Q}{4\pi\varepsilon b^2}=-\frac{Q}{4\pi b^2} \]
\[ Q_b=\int_{\theta=0}^{\pi}\int_{\phi=0}^{2\pi}-\frac{Q}{4\pi b^2}\,b^2\sin\theta\,d\phi\,d\theta=-Q \]
\[ Q_c=-Q_b=Q \]

\noindent c)\ \ $r>c$ iken elektrik alan nedir??
\[ \oint_S \vec D\cdot d\vec s=Q+Q_b+Q_c=Q \]
\[ \vec E=\frac{Q}{4\pi\varepsilon r^2}\,\vec e_r \]

\noindent d)\ \ çepere dokundurup uzaklaştırsak

\hspace{1cm}kürecik nötrleşir

\end{document}
